\section{Wprowadzenie}
\label{sec:intro}

Grafika komputerowa 3D znajduje zastosowanie w niemal każdej dziedzinie, od przemysłu rozrywkowego, przez architekturę, aż do medycyny, gdzie pomaga wizualizować obiekty trójwymiarowe. Modelowanie geometryczne jest jednym z jej współczesnych fundamentów. Sposób w jaki komputer reprezentuje i przetwarza kształty trójwymiarowe ma bezpośredni wpływ na jakość obrazu, wydajność obliczeniową oraz możliwości edycji i animacji obiektów.

Celem pracy jest przegląd i analiza metod stosowanych do wspomagania modelowania geometrycznego. Omówiono trzy główne obszary, które zostały wybrane do analizy, ponieważ reprezentują komplementarne elementy modelowania geometrycznego: sposób przechowywania geometrii, metody jej generowania i modyfikacji oraz sposoby interakcji.

Rozdział 2 przedstawia sposoby reprezentacji geometrii. Rozdział 3 poświęcony jest proceduralnemu generowaniu geometrii. Rozdział 4 opisuje algorytmy stosowane na modelach.

Praca opiera się na przeglądzie literatury naukowej z zakresu grafiki komputerowej, obejmującej książki jak i artykuły konferencyjne publikowane na SIGGRAPH\footnote{Special Interest Group on Computer Graphics and Interactive Techniques - największa coroczna konferencja naukowa z zakresu grafiki komputerowej} oraz w czasopismach naukowych.